\documentclass[11pt]{article}
\usepackage[margin=1in]{geometry}
\usepackage{amsmath,amssymb,amsthm,mathtools}
\usepackage{bm}
\usepackage{hyperref}
\usepackage{enumitem}
\usepackage{booktabs}
\usepackage{array}
\usepackage{longtable}
\usepackage{multirow}
\usepackage{graphicx}
\usepackage{setspace}
\usepackage{microtype}
\usepackage[T1]{fontenc}
\usepackage[utf8]{inputenc}

\title{From the Growing Ledger to Known Physics: \\ Gravity, Matter, and Coupling Structure in Relational Actualism}
\author{Joshua Sandeman}
\date{\today}

\begin{document}
\maketitle

\begin{abstract}
Relational Actualism (RA) proposes that physical reality is fundamentally composed not of objects in a pre-given spacetime, but of irreversible actualization events whose accumulation forms a growing causal ledger. A companion theory-definition paper presents the process ontology, the closure architecture by which that ontology becomes a specific implemented $d=4$ theory, and the discrete substrate on which the framework rests. The present paper develops the bridge from implemented $d=4$ RA to known physics.

It presents four principal bridge domains: the gravity bridge, the matter bridge, the gauge and interaction bridge, and the motif-renewal lifetime bridge. A central claim is that the implemented $d=4$ theory now supports a significantly stronger bridge architecture than earlier versions of the corpus. The local BDG coefficients are fixed by dimensional closure; the actualization threshold fixes a candidate discrimination length; that discrimination scale fixes a candidate hadronic density scale; the hadronic density scale fixes the raw direct-exit geometry of openly destabilizable motifs; and additional stabilization layers stratify the longer lifetime hierarchy.

Throughout, the paper adopts explicit status discipline. Lean-verified, computation-verified, derived, interpretive, and open-target components are separated rather than allowed to slide rhetorically into one another.
\end{abstract}

\section*{Front-matter status box}
\noindent\textbf{This paper does:} present the strongest bridge results from implemented Relational Actualism to known physics; separate Lean-verified, computation-verified, derived, and interpretive layers; and show how gravity, matter structure, interaction hierarchy, and particle lifetimes arise from the implemented $d=4$ theory.

\noindent\textbf{This paper does not:} redefine the kernel from scratch, reargue the full ontology, or present the entire prediction program.

\section{Introduction}
Paper~I (\cite{RAKernelPaper}) defines Relational Actualism as a process ontology of irreversible actualization and presents the closure architecture by which that ontology becomes a specific implemented $d=4$ theory. The purpose of the present paper is different. It asks what known physics follows from that implemented structure, and with what epistemic status.

The paper has four principal bridge domains. First, it presents the gravity bridge: the route from local BDG action, ledger conservation, and continuum closure to Einstein dynamics. Second, it presents the matter bridge: the route from finite topology closure to particle-class structure. Third, it presents the gauge and coupling bridge: the route from depth/sign structure to effective interaction hierarchy. Fourth, it presents the decay and persistence bridge: the route from motif renewal, density-controlled exit geometry, and $\sigma$-filtered path structure to the observed lifetime hierarchy.

This last bridge has become substantially stronger in recent work. The decay program no longer looks like a collection of phenomenological adjustments surrounding a discrete core. It now has a discernible architecture. $d=4$ closure fixes the integer substrate; that substrate fixes the selectivity threshold structure; the threshold structure fixes a candidate hadronic density scale; that density fixes the raw direct-exit geometry of openly destabilizable motifs; and additional stabilization layers then stratify the longer lifetime hierarchy through filtered exit paths.

\section{The gravity bridge}
\textbf{Main status of this section:} \textbf{mixed (LV + published continuum inputs + DR + IN)}.

The gravity bridge is the strongest currently available route from implemented RA to familiar continuum physics. Its central claim is that the local BDG action, together with event-level conservation and dense-regime continuum closure, yields Einstein dynamics in the appropriate limit.

The first route proceeds through Lovelock uniqueness. The second proceeds more directly through the discrete substrate:
\begin{enumerate}[label=\arabic*.]
    \item the local action coefficients are fixed by dimensional closure in $d=4$;
    \item the exact second-order cancellation condition $\sum_k c_k = 0$ ensures the d'Alembertian limit;
    \item the Benincasa--Dowker continuum result then links the local action to Einstein--Hilbert structure.
\end{enumerate}

\section{Stable topology and the finite local matter universe}
\textbf{Main status of this section:} \textbf{LV + CV}.

The matter bridge begins not with particle names, but with topology closure. The strongest theorem-level statement in this domain is that the local stable topology universe in implemented $d=4$ RA closes into a finite set of stable classes. The relevant closure theorem, referred to in the internal architecture as L11, states that the stable local motif universe closes into exactly five classes under exhaustive single-step extension analysis.

\section{From topology classes to particle classes}
\textbf{Main status of this section:} \textbf{mixed (LV/CV + IN)}.

The finite topology universe does not by itself name quarks, leptons, gauge bosons, or Higgs-like motifs. The Standard Model reading is an interpretive bridge built on top of the topology closure theorem. The strongest careful statement here is:
\begin{quote}
The topology classification is derived. The Standard Model identification is a systematic interpretive mapping built on that classification.
\end{quote}

\section{Gauge structure and force differentiation}
\textbf{Main status of this section:} \textbf{mixed (CV + DR + IN)}.

The BDG vector in $d=4$ exhibits a nontrivial alternating-sign structure:
\[
(1,-1,9,-16,8).
\]
This structure already supports a native distinction between destabilizing and stabilizing channels. The force-differentiation program suggests that the relevant effective channel weights take the form
\[
W_k(\mu) \sim |c_k|\,\frac{\mu^k}{k!},
\]
so that all channels are comparable at $\mu\sim 1$ while differentiating factorially at lower densities.

\section{Motif renewal, selection rules, and particle lifetimes}
\textbf{Main status of this section:} \textbf{mixed (CV + DR + SYN + OPEN)}.

In Relational Actualism, a particle is a persistent causal motif represented by a fuller motif state
\[
M=(T,N,C,\sigma),
\]
where $T$ is topology class, $N$ is the local BDG depth profile, $C$ is confinement/reset memory, and $\sigma$ is the internal label structure relevant to accessible renewal and exit channels.

Lifetime is modeled as exit geometry in topology space:
\[
\tau_P^{(\mathrm{steps})}
=
\mathbb E\!\left[\inf\{t>0:M_t\notin [M]_P\}\right].
\]

The direct same-force decay program becomes substantially stronger if its density scale is not fitted but derived. The candidate closure is
\[
l_{\mathrm{RA}}:=\sqrt{4\Delta S^*}, \qquad
\mu_{\mathrm{QCD}}:=e^{l_{\mathrm{RA}}}.
\]
With $\Delta S^*=0.60069$ nats, one obtains $l_{\mathrm{RA}}\approx 1.550$ and therefore $\mu_{\mathrm{QCD}}\approx 4.7119$.

The direct strong-decay sector is the cleanest place where the derived density should first appear. These are Type~I motifs: direct-exit motifs with $\Sigma_{\mathrm{eff}}\sim 1$.

The observed hierarchy of longer lifetimes is interpreted as a hierarchy of stabilization layers. Recent empirical synthesis suggests a five-layer classification:
\[
\text{Type I} > \text{Type II} > \text{Type III} > \text{Type IV} > \text{Type V}
\]
in accessible exit strength, and therefore the reverse in lifetime.

Once a motif's shortest direct exit is blocked, the relevant object becomes a pathwise exit weight:
\[
P(\gamma\mid M_0,\mu)
=
\prod_{j=0}^{n-1}
\Sigma_{k_j}(M_j)\,
A_{k_j}(M_j)\,
G_{k_j}(\mu)\,
B_{k_j}(M_j\to M_{j+1}),
\]
and the total exit probability is
\[
P_{\mathrm{exit}}(M\mid\mu)
=
\sum_{\gamma\in\Gamma_{\mathrm{exit}}(M)}
\left[
\prod_{j}
\Sigma_{k_j}(M_j)\,
A_{k_j}(M_j)\,
G_{k_j}(\mu)\,
B_{k_j}(M_j\to M_{j+1})
\right]V(\gamma).
\]

\section{Conclusion}
The bridge program is no longer a collection of unrelated claims. It has an internal order. Gravity emerges from the discrete action and its continuum limit. Matter emerges from finite topology closure. Interaction hierarchy emerges from depth-channel structure. Persistence and decay emerge from renewal, filtered exit geometry, and stabilization layers.

\bibliographystyle{plain}
\bibliography{ra_suite_master_references}
\end{document}
